\chapter{Context: reinforcement learning and the sixteen flaws}\label{chap:02-contexte}

This chapter sets out the notions the rest of the report depends on: the formal framework of
reinforcement learning (Section~\ref{sec:cadre}), the four programs of our experiments and the TRPO
algorithm (Section~\ref{sec:ctx-programmes}), the simulated tasks (Section~\ref{sec:ctx-taches}), what
a random seed fixes and how an agent is evaluated (Sections~\ref{sec:ctx-graines}
and~\ref{sec:ctx-evaluation}). It then inventories the sixteen flaws Joelle Pineau showed in 2018
(Section~\ref{sec:failles}), places our work among related work (Section~\ref{sec:voisins}) and fixes
the vocabulary used for language models (Section~\ref{sec:vocabulaire}).

\section{The formal framework}\label{sec:cadre}

\paragraph{The Markov decision process}
Reinforcement learning (RL) studies a program, the agent, that acts repeatedly in an environment. The
standard framework is the Markov decision process \parencite{puterman1994mdp,sutton2018rl}, a 5-tuple
\begin{equation}
  \mathcal{M} = (\mathcal{S}, \mathcal{A}, P, r, \gamma),
\end{equation}
where $\mathcal{S}$ is the set of states, $\mathcal{A}$ the set of actions, $P(s' \mid s, a)$ the
probability of moving to state $s'$ after action $a$ in state $s$, $r(s, a)$ the reward received and
$\gamma \in [0, 1)$ the discount factor, which weights a reward less the later it comes. A distribution $\rho_0$ draws
the initial state. In the tasks of this report, $\mathcal{S}$ and $\mathcal{A}$ are subsets of
$\mathbb{R}^d$: the state describes the positions and velocities of the joints of a simulated body,
the action the torques applied to its motors.

\paragraph{Policy, episode and return}
A policy $\pi_\theta(a \mid s)$ gives the distribution of the action in each state. In deep RL, the
policy is a neural network with parameters $\theta$. An episode is a sequence
$\tau = (s_0, a_0, r_0, s_1, a_1, r_1, \dots, s_T)$ obtained by drawing $s_0 \sim \rho_0$, then
$a_t \sim \pi_\theta(\cdot \mid s_t)$ and $s_{t+1} \sim P(\cdot \mid s_t, a_t)$, until a terminal
state or a maximum duration. The discounted return of an episode and the learning objective are
\begin{equation}
  G(\tau) = \sum_{t=0}^{T-1} \gamma^t \, r(s_t, a_t),
  \qquad
  J(\theta) = \mathbb{E}_{\tau \sim \pi_\theta}\bigl[ G(\tau) \bigr].
\end{equation}
Learning seeks the $\theta$ that maximises $J$. The score reported in papers and in our experiments is
almost always the undiscounted return $\sum_t r(s_t, a_t)$ of an episode, averaged over several
evaluation episodes.

\paragraph{Value functions}
The value of a state and the value of an action under the policy $\pi$ are
\begin{equation}
  V^\pi(s) = \mathbb{E}_\pi\Bigl[ \sum_{k \geq 0} \gamma^k r_{t+k} \Bigm| s_t = s \Bigr],
  \qquad
  Q^\pi(s, a) = \mathbb{E}_\pi\Bigl[ \sum_{k \geq 0} \gamma^k r_{t+k} \Bigm| s_t = s, a_t = a \Bigr].
\end{equation}
The advantage $A^\pi(s, a) = Q^\pi(s, a) - V^\pi(s)$ measures how much more an action yields than the
average action of the policy in that state. An actor-critic algorithm learns both a policy (the actor)
and an estimate of a value function (the critic). The critic's estimate reduces the variance of the
actor's learning signal.

\paragraph{Policy gradient}
For a policy differentiable in $\theta$, the policy gradient theorem \parencite{sutton2018rl} gives
\begin{equation}
  \nabla_\theta J(\theta) = \mathbb{E}_{\pi_\theta}\bigl[ \nabla_\theta \log \pi_\theta(a_t \mid s_t)\, A^{\pi_\theta}(s_t, a_t) \bigr].
\end{equation}
So-called \emph{on-policy} methods, such as TRPO and PPO, estimate this expectation with episodes drawn
from the current policy, then discard these episodes after the update. \emph{Off-policy} methods, such
as SAC and CrossQ, keep past transitions in a replay buffer and learn a critic $Q$ from them. They reuse
each transition several times. For the same number of environment steps, off-policy methods generally
learn faster, at the price of heavier computation per step.

\paragraph{Hyperparameters}
The parameters $\theta$ are learned. Other settings stay fixed by the experimenter: the learning rate,
$\gamma$, the size of the networks, the batch size, the number of steps collected before each update.
These are the hyperparameters. Flaws F4 to F7 of Section~\ref{sec:failles} concern them.

\section{The four programs and TRPO}\label{sec:ctx-programmes}

Our experiments use four programs: four specific implementations of an algorithm, each in a given
library and version. The word \enquote{algorithm} denotes the method described in a paper. The word
\enquote{program} denotes the code that runs it. Two programs of the same algorithm sometimes give
different results, which is flaw F14. The pinned versions of each library are given in
Table~\ref{tab:env}. Chapter~\ref{chap:05-experiences-methode} details the budgets and the evaluation.

\paragraph{PPO, in PyTorch and in JAX}
PPO (\emph{Proximal Policy Optimization}) \parencite{schulman2017ppo} is an \emph{on-policy}
actor-critic algorithm. With the probability ratio
$\rho_t(\theta) = \pi_\theta(a_t \mid s_t) / \pi_{\theta_\text{old}}(a_t \mid s_t)$ and an estimator
$\hat{A}_t$ of the advantage, PPO maximises the clipped objective
\begin{equation}
  L^{\text{CLIP}}(\theta) = \hat{\mathbb{E}}_t\Bigl[ \min\bigl( \rho_t(\theta)\, \hat{A}_t,\;
  \operatorname{clip}\bigl(\rho_t(\theta), 1 - \epsilon, 1 + \epsilon\bigr)\, \hat{A}_t \bigr) \Bigr],
\end{equation}
to which a critic loss and an entropy bonus are added. Clipping removes the incentive for an update to
move the new policy far from the old one. The program \algo{PPO PyTorch} is the one in
Stable-Baselines3 \parencite{raffin2021sb3}. The program \algo{PPO JAX} is the one in sbx, the JAX
rewrite of Stable-Baselines \parencite{raffin2023sbx,bradbury2018jax}. Both programs receive the
same algorithm hyperparameters and collect experience in batches of 2,048 steps. Their default
initialisation differs. Stable-Baselines3 turns on its \file{ortho\_init} option: the hidden layers
receive an orthogonal initialisation with gain $\sqrt{2}$, the action output a gain of 0.01 and the
value output a gain of 1 (\file{stable\_baselines3/common/policies.py}). The sbx library turns it off by
default and keeps Flax's initialisation. When turned on, sbx's option touches only the action output
layer (\file{sbx/ppo/policies.py}). Experiment E3 compares these two programs
(Chapter~\ref{chap:06-experiences-resultats}). PPO was the standard baseline of 2018. The protocol keeps it because it remains a
common baseline in 2026 and because GRPO, the algorithm of our language-model strand
(Chapter~\ref{chap:07-modeles-de-langue}), derives from it.

\paragraph{SAC}
SAC (\emph{Soft Actor-Critic}) \parencite{haarnoja2018sac} is an \emph{off-policy} actor-critic
algorithm. It maximises the return augmented by the entropy of the policy,
\begin{equation}
  J_{\text{SAC}}(\pi) = \sum_t \mathbb{E}_{(s_t, a_t) \sim \pi}\bigl[ r(s_t, a_t) + \alpha\, \mathcal{H}\bigl(\pi(\cdot \mid s_t)\bigr) \bigr],
\end{equation}
where the coefficient $\alpha$ sets the weight given to exploration and adapts during training. Two $Q$ critics learned
in parallel limit the overestimation of values. Delayed copies of these critics, the target networks,
stabilise the learning targets. The \algo{SAC} program of our experiments is the one in sbx. SAC is a
common \emph{off-policy} baseline in continuous control.

\paragraph{CrossQ}
CrossQ \parencite{bhatt2024crossq} descends from SAC. It removes the target networks and normalises
the activations by batch. The pairs $(s, a)$ and $(s', a')$ go through the critic together, so
that the normalisation sees both distributions. Its authors report a sample efficiency, the score
reached per environment step, at least equal to that of methods far more expensive in computation. The
\algo{CrossQ} program of our experiments is the one in sbx. It uses batch renormalisation
\parencite{ioffe2017batchrenorm}, in the actor as in the critic, with a momentum of 0.99. For the first
100,000 updates of each network, batch renormalisation behaves as plain batch normalisation. Measured
against our budget of steps, this threshold keeps the critic in that regime for the first half of
training. The actor, updated once in three, never reaches the threshold
(Section~\ref{sec:programmes}). Its Adam optimiser takes $\beta_1 = 0.5$ instead of 0.9
(\file{sbx/crossq/policies.py}). Table~\ref{tab:hyperparametres} gives its other settings.

\paragraph{TRPO}
TRPO (\emph{Trust Region Policy Optimization}) \parencite{schulman2015trpo} is the algorithm of the most
famous figure of 2018 (flaw F11). It maximises a surrogate objective under a Kullback-Leibler
divergence constraint,
\begin{equation}
  \max_\theta\; \hat{\mathbb{E}}_t\bigl[ \rho_t(\theta)\, \hat{A}_t \bigr]
  \quad \text{subject to} \quad
  \hat{\mathbb{E}}_t\bigl[ \mathrm{KL}\bigl( \pi_{\theta_\text{old}}(\cdot \mid s_t) \,\big\|\, \pi_\theta(\cdot \mid s_t) \bigr) \bigr] \leq \delta,
\end{equation}
a problem solved by conjugate gradient followed by a line search. PPO is a simplification of TRPO. We
use TRPO only in experiment B1, with the sb3-contrib~2.9.2 program and its default values, without any
tuning (Chapter~\ref{chap:05-experiences-methode}).

\section{The MuJoCo tasks}\label{sec:ctx-taches}

Henderson and co-authors compared their algorithms on tasks of the MuJoCo simulator (\emph{Multi-Joint
dynamics with Contact}) \parencite{todorov2012mujoco,henderson2018matters}. We take the same tasks in
their 2026 version, the one in Gymnasium \parencite{towers2024gymnasium}. Table~\ref{tab:taches} gives
the dimensions of the state and action spaces, read in the pinned environment of our experiments. Each
episode lasts at most 1,000 steps. \task{Hopper} and \task{Walker2d} stop earlier when the body
falls.

\begin{table}[tbp]
\centering\small
\caption{The four tasks of our experiments, in Gymnasium~1.1.1 and MuJoCo. The dimensions and the
maximum length of an episode are read in the environment of our experiments (\file{sim/env}).}
\label{tab:taches}
\begin{tabular}{@{}l r r >{\raggedright\arraybackslash}p{7.4cm}@{}}
\toprule
Task & state & action & Simulated body and reward \\
\midrule
\task{HalfCheetah-v5} & 17 & 6 & articulated cheetah in two dimensions. Reward: forward velocity, minus a control cost \\
\task{Hopper-v5} & 11 & 3 & one-legged hopper. Reward: forward velocity, plus a bonus while the body stays upright \\
\task{Walker2d-v5} & 17 & 6 & two-legged walker. Reward: the same terms as \task{Hopper} \\
\task{Swimmer-v5} & 8 & 2 & three-segment swimmer in a viscous fluid. Reward: forward velocity \\
\bottomrule
\end{tabular}
\end{table}

Pineau judged these benchmarks \enquote{ridiculously simple}: states described by few numbers,
deterministic transitions, a world weighing less than 100~kilobytes \parencite{pineau2018neurips}. This
criticism grounds flaw F16. We keep these tasks, since the aim is to rerun the experiments of 2018 as
they were run.

\section{What a random seed changes}\label{sec:ctx-graines}

\paragraph{Definition}
A training run draws many numbers at random. A computer produces this randomness from a starting
integer, the random seed. In our programs, the random seed is passed to the library (\file{seed=SEED}
in \file{sim/run.py}). The library derives from it the initialisation of the network weights, the
draw of exploration actions, the order of the mini-batches and the sequence of initial states of the
training environment. In this report, a \emph{run} denotes a complete training with a given random
seed.

\paragraph{What stays fixed}
Same program, same random seed, same library versions and same machine: the protocol expects the same
result. Our logs confirm it for \algo{PPO PyTorch} and \algo{PPO JAX}: a run stopped then restarted recovers its evaluation
points identically. For \algo{SAC} and \algo{CrossQ}, three training runs redone unchanged recover their
evaluation points to the byte (\cref{sec:outils}).

\paragraph{What changes}
Changing the random seed changes the initialisation of the weights, the exploration, the order of the
mini-batches and the initial states. The final score of a given program then varies a great deal. Over our thirty runs of
\algo{PPO PyTorch} on \task{HalfCheetah}, the standard deviation is 34\% of the mean
(Chapter~\ref{chap:06-experiences-resultats}).

\paragraph{The statistical unit is the run}
The random seed is a source of variance, in the same way as the draw of a sample in a survey. A
comparison between two programs therefore rests on several independent runs per program, with an
uncertainty computed over these runs. With one run per program, the comparison measures the luck of
the draw. An episode, or a point on a curve, does not count as an independent observation, since it
belongs to one run. Flaw F11 and experiment B1 show what happens when the episode is taken as the unit
(Chapter~\ref{chap:06-experiences-resultats}). For this reason, criterion X8 of the survey asks whether
papers compute their uncertainty over independent runs (Section~\ref{sec:sm-x8}).

\section{Evaluating an agent}\label{sec:ctx-evaluation}

\paragraph{The score of a run}
Evaluating an agent means playing episodes with its policy, without learning, then averaging their
returns. Three choices remain open. A paper that leaves them unstated shows flaw F8.
\begin{itemize}
  \item The measurement point: the score at the end of training, the best score reached along the
  curve, or the mean of the curve, which approximates the area under the curve divided by its length.
  \item The policy mode: stochastic as in training, or deterministic (the most probable action).
  \item The environment: the training one, or held-out conditions, such as other environment random
  seeds or modified physics (flaw F15).
\end{itemize}
Our programs are evaluated every 50,000 steps on 10 episodes, in deterministic mode, with
environment random seeds $10{,}000 + i$, where $i$ numbers the evaluation episodes. These random
seeds never serve in training (Chapter~\ref{chap:05-experiences-methode}).

\paragraph{Aggregating runs}
With $n$ runs of scores $x_{(1)} \leq \dots \leq x_{(n)}$, the mean and the median react differently to
extreme runs. \textcite{agarwal2021precipice} recommend the interquartile mean (IQM), which discards
the lowest quarter and the highest quarter of the runs, then averages the rest:
\begin{equation}
  \mathrm{IQM}(x) = \frac{1}{n - 2k} \sum_{i = k + 1}^{n - k} x_{(i)}, \qquad k = \lfloor n/4 \rfloor .
\end{equation}
The IQM is the 25\% trimmed mean of \file{scipy.stats.trim\_mean}, which the rliable library calls:
with ten runs, the IQM is the mean of the six middle runs. Its confidence interval comes from a
stratified bootstrap. The runs are resampled with replacement within each task, the IQM is
recomputed, then the 2.5\% and 97.5\% quantiles of the 2,000 replicates give the interval
(\file{src/pineau\_rl/bootstrap.py}).

\paragraph{Comparing two programs}
The probability of improvement of $X$ over $Y$ is the probability that a run of $X$ beats an
independently drawn run of $Y$. With $n_X$ runs of $X$ and $n_Y$ runs of $Y$, it is estimated by
\begin{equation}
  P(X > Y) = \frac{1}{n_X n_Y} \sum_{i=1}^{n_X} \sum_{j=1}^{n_Y} \Bigl( \mathbb{1}[x_i > y_j] + \tfrac{1}{2}\, \mathbb{1}[x_i = y_j] \Bigr).
\end{equation}
This estimate equals the Mann-Whitney $U$ statistic divided by $n_X n_Y$. At 0.5, a run of $X$ beats a
run of $Y$ half the time. Our experiments declare a gap settled when the 95\% bootstrap interval of
$P(X > Y)$ excludes 0.5, or when the intervals of the two IQMs do not overlap: this is the
settled-winner rule. On groups of runs of one program, where no gap exists, it concludes wrongly slightly more than once in twenty with
ten runs per group, nearly once in ten with five (Section~\ref{sec:regle}). Welch's $t$-test \parencite{welch1947generalization}, which compares two
means without assuming equal variances, serves only in experiments E1 and B1, since the 2018 figure
used it.

\section{The sixteen flaws shown in 2018}\label{sec:failles}

In her two talks of 2018, at ICLR and then at NeurIPS \parencite{pineau2018iclr,pineau2018neurips},
Pineau shows each flaw with an example, often a figure taken from her team's paper signed by Peter
Henderson and co-authors \parencite{henderson2018matters}. We recorded sixteen flaws, in the order of
the two talks, then sorted them into four families. Table~\ref{tab:failles} lists them with the slide
or page that shows each one and the instrument of this report that measures it. The labels F1 to F16
serve only to find a flaw in the tables. Appendix~\ref{chap:A-grille} gives the decision rules of
each reading criterion.

\begin{small}
\begin{longtable}{@{}>{\raggedright\arraybackslash}p{0.6cm} >{\raggedright\arraybackslash}p{5.3cm} >{\raggedright\arraybackslash}p{3.6cm} >{\raggedright\arraybackslash}p{4.4cm}@{}}
\caption{The sixteen flaws Joelle Pineau showed in 2018 in deep reinforcement learning. For each one:
the slide of her ICLR talk or the page of her NeurIPS document that shows it, then the reading
criteria (Q, X, Y, Chapter~\ref{chap:03-sondage-methode}) and the experiments (E, B,
Chapter~\ref{chap:05-experiences-methode}) that measure it in this report.}\label{tab:failles}\\
\toprule
& Flaw & Where Pineau shows it & What measures it here \\
\midrule
\endfirsthead
\toprule
& Flaw & Where Pineau shows it & What measures it here \\
\midrule
\endhead
\bottomrule
\endfoot
\multicolumn{4}{@{}l}{\emph{The materials: what rerunning takes}} \\
F1 & The paper published without its code & slide 6 (ICLR), page 29 (NeurIPS) & criterion Q5 \\
F2 & No plan for rerunning the experiment, as with AlphaGo & slides 9 to 12 (ICLR) & criteria Q3 and X7 \\
F3 & Hardware and software, which belong to the result & pages 34 and 35 (NeurIPS) & criteria X7 and Y10 \\
\midrule
\multicolumn{4}{@{}l}{\emph{Tuning: hyperparameters and baselines}} \\
F4 & The method's hyperparameters, given in part & page 29 (NeurIPS) & criterion Q3 \\
F5 & How the hyperparameters were chosen, left unsaid & page 29 (NeurIPS) & criterion Q6 \\
F6 & Settings that change everything and interact & slides 24 and 25 (ICLR) & criterion Y9 \\
F7 & Baselines tuned less carefully than the method & slide 25 (ICLR), page 29 (NeurIPS) & criteria Q4 and X4, experiment E6 \\
\midrule
\multicolumn{4}{@{}l}{\emph{Random seeds and statistics}} \\
F8 & A performance measure never defined & slides 26 and 27 (ICLR) & criteria Y1 and Y2, experiment B3 \\
F9 & Too few runs: three or five, never ten in her table & slide 28 (ICLR) & criterion X1, experiment E5 \\
F10 & Only the best runs out of $N$ published & slides 29 and 30 (ICLR) & criterion X3, experiments E1 and E4 \\
F11 & Same code, another random seed: two \enquote{algorithms} & slide 31 (ICLR), page 22 (NeurIPS) & experiments E1 and B1 \\
F12 & No statistical test & page 29 (NeurIPS) & criteria Q8, X2, Y11 and Y12 \\
\midrule
\multicolumn{4}{@{}l}{\emph{The reach of the conclusions}} \\
F13 & The best algorithm changes with the task & slides 20 and 21 (ICLR) & criteria Y3 and Y4, experiment E2 \\
F14 & Two implementations of the same algorithm that diverge & slides 22 and 23 (ICLR) & criterion Y5, experiment E3 \\
F15 & Evaluation in the training conditions & pages 36 to 51 (NeurIPS) & criteria Q7 and Y8, experiment E7 \\
F16 & Benchmarks too simple, easy to learn by heart & pages 43 to 50 (NeurIPS) & criteria X6 and Y7 \\
\end{longtable}
\end{small}

\subsection{The materials: what rerunning takes}

\paragraph{F1, the paper published without its code}
An RL paper describes its method in a few pages. The program that produced the results runs to
thousands of lines, full of choices the prose never mentions. Pineau takes up the line of
\textcite{buckheit1995wavelab} quoted in Chapter~\ref{chap:01-introduction}, in which the two authors
condensed an idea of Jon Claerbout: the paper is the advertising, the research is the code and the instructions that
produced the figures. Of the 50 RL papers from 2018 she read, 55\% link to their code.

\paragraph{F2, no plan for rerunning the experiment}
Code falls short without the exact settings and an idea of the compute required. Pineau's example is
AlphaGo \parencite{silver2016alphago}. The paper is detailed, with a long appendix. Her slide still
asks whether that is enough: a rerun would also need the expert games, the code, the hyperparameters,
the computers and the evaluation protocol. A discovery that no outside lab has the means to redo stays
a claim.

\paragraph{F3, hardware and software}
The same code launched on two machines sometimes gives two results. Graphics cards round computations
differently from processors. A new version of a library changes a default value. Pineau's checklist
asks for \enquote{the computing infrastructure used} and for the code \enquote{including all
dependencies}.

\subsection{Tuning: hyperparameters and baselines}

\paragraph{F4, the method's hyperparameters, given in part}
One setting left out blocks reproduction: the reader takes the default value of their library and
gets something else. Pineau counts 90\% of 2018 papers giving all the hyperparameters of their method.

\paragraph{F5, how the hyperparameters were chosen, left unsaid}
A final value says nothing of how it was found. A team that tried a hundred combinations and kept the
best one has optimised on its benchmark. The reader needs the search budget. Only 20\% of 2018 papers
state it.

\paragraph{F6, settings that change everything}
Henderson and co-authors change one setting at a time. For PPO on \task{HalfCheetah}, going from a
network of two layers of 64 neurons to a network of 400 then 300 neurons drops the score from
2,201 to $-1{,}180$. The settings interact: the best network for one algorithm stops being the
best for another. \textcite{engstrom2020implementation} and then \textcite{andrychowicz2021matters}
confirmed this sensitivity at large scale.

\paragraph{F7, baselines tuned less carefully than the method}
The author of a new method tunes it with care. The same author often takes the baseline with its
default values. The comparison then pits a tuned method against an untuned one. Pineau counts 60\% of
2018 papers giving the hyperparameters of their baselines. \textcite{eimer2023hyperparameters} call for
the same tuning budget for both.

\subsection{Random seeds and statistics}

\paragraph{F8, a performance measure never defined}
The final score, the mean over training and the best score reached do not always rank methods in the
same order (Section~\ref{sec:ctx-evaluation}). A paper that states neither which measure it uses nor
what its error bars represent compares with no other paper.

\paragraph{F9, too few runs}
Pineau shows a table of papers from 2015 to 2017: they launch three or five random seeds, none of them
ten. With three runs, a mean depends heavily on the luck of one run.
\textcite{colas2018seeds} compute the number of random seeds needed to detect a given gap.

\paragraph{F10, the best runs out of $N$}
A team launches ten runs, keeps the best ones and publishes their mean. Every number is true, yet the
published mean exceeds what the method delivers on average. The spread also looks smaller than it is.
Pineau illustrates the practice with ten curves of which she keeps the best three.

\paragraph{F11, same code, another random seed}
Henderson's team trains TRPO ten times on \task{HalfCheetah}, changing only the random seed.
They form two groups of five and plot the two mean curves, which stand clearly apart. A test declares
the curves different ($t = -9.09$, $p = 0.0016$). On stage, Pineau specifies that the groups are not
sorted. Chapters~\ref{chap:05-experiences-methode} and~\ref{chap:06-experiences-resultats} return at
length to this number (experiments E1 and B1).

\paragraph{F12, no statistical test}
When two methods differ by a few points, a test or an interval tells whether the gap exceeds what the
luck of the random seeds produces by itself. In Pineau's reading, 5\% of 2018 papers apply a test to
their comparisons.

\subsection{The reach of the conclusions}

\paragraph{F13, the best algorithm changes with the task}
Henderson's paper ranks four algorithms on four tasks. The algorithm that leads on
\task{HalfCheetah} stops leading on \task{Hopper} or \task{Swimmer}. A general conclusion drawn from
one or two tasks generalises a special case.

\paragraph{F14, two implementations that diverge}
Three implementations of TRPO, with the same network, finish on \task{HalfCheetah} at 1,939, 834
and 205. Comparing a method with \enquote{TRPO} without saying which implementation was taken amounts to
comparing it with an anonymous program.

\paragraph{F15, evaluation in the training conditions}
In supervised learning, a model is evaluated on data it has never seen. In RL, the agent is often
evaluated in the same world it trained in. Pineau shows a sentence she puts up for debate, according to
which RL would be the only case of machine learning where testing on the training set is acceptable.
She proposes separating the training and test random seeds, then the tasks
\parencite{zhang2018dissection,zhang2018natural}. Pineau finds only 10\% of 2018 papers that
evaluate on held-out conditions.

\paragraph{F16, benchmarks too simple}
A large enough agent memorises the right sequence of actions on a deterministic benchmark without
understanding anything about the world \parencite{machado2018revisiting}. Pineau also asks for the
version of each benchmark, so that a result can be rerun in the same world.

\paragraph{The remedies of 2018}
Against the sixteen flaws, Pineau first sets a reproducibility checklist (pages 31 to 33 at NeurIPS).
For each algorithm, it asks for a clear description, a complexity analysis, the link to the code and
its dependencies. For each theoretical result, it asks for the statement, the assumptions and the
proof. For each figure, it asks for the data collection, the split between training, validation and
test, the excluded data, the range and selection method of the hyperparameters, the exact number of
runs, the definition of the statistics and error bars, the hardware used. Pineau adds the ICLR 2018
reproducibility challenge (slides 37 to 43 at ICLR), presented in Chapter~\ref{chap:01-introduction}.
NeurIPS adopted the checklist in 2019 \parencite{pineau2021improving}.

\section{Related work}\label{sec:voisins}

At the start of the project, we searched arXiv, OpenReview and the conference proceedings. To our knowledge, no
work reapplies Pineau's grid to recent RL papers while calibrating it on 2018. None reruns
Henderson's protocol with the algorithms of 2026 under a protocol written in advance. Related
work falls into five groups.

\paragraph{The 2018 experiment}
\textcite{henderson2018matters} ground most of flaws F6 to F14: sensitivity to hyperparameters, to
implementations, to random seeds, to tasks. \textcite{islam2017reproducibility} preceded them on the
same tasks. \textcite{lynnerup2020survey} rerun algorithms on a real robot arm.
\textcite{bjorck2022variance} show, for control from images, that part of the variance between random
seeds comes from runs that fail early through numerical instability.

\paragraph{Statistical method}
\textcite{colas2018seeds} run a power analysis of the number of random seeds.
\textcite{jordan2020evaluating} discuss the choice of the performance measure, then
\textcite{jordan2024position} the cost of a rigorous comparison. \textcite{agarwal2021precipice}
contribute the IQM, stratified bootstrap intervals, performance profiles and the probability of
improvement, gathered in the rliable library. They show that several improvements claimed on Atari
games dissolve into the uncertainty once the intervals are computed. \textcite{patterson2024empirical}
turn these results into a guide. \textcite{eimer2023hyperparameters} call for an equal tuning budget
for baselines. These works say what to do. Our survey measures whether papers do it.

\paragraph{Checklists and reproducibility challenges}
\textcite{pineau2021improving} describe the NeurIPS 2019 reproducibility programme: the mandatory
checklist, the reproducibility challenge and the code submission policy. The share of accepted papers
with code rises from under 50\% to 74.4\%. \textcite{magnusson2023nlpchecklist} measure the effect of
the checklist of natural language processing conferences.

\paragraph{Audits of the literature}
\textcite{gundersen2018state} score 400 papers from the IJCAI and AAAI conferences on how they document
methods, data and experiments. \textcite{coakley2026shift} track 56,800 papers from five
artificial intelligence conferences from 2014 to 2024: the share that releases code and data rises
from 11\% to 64\%. The rise begins before the checklists. In RL, the closest audit is the one by
\textcite{gorsane2022marl}: 75 papers on cooperative multi-agent RL, from 2016 to 2022, of which about
a third quantify no uncertainty. These audits leave aside single-agent RL and the 2018 grid.

\paragraph{Language models}
\textcite{hochlehnert2025sober} re-evaluate reasoning models and find standard deviations of 5 to 15
percentage points from one sampling random seed to the next on small benchmarks such as AIME~2024.
\textcite{chandak2025incorrect} report, in a blog post that was not peer-reviewed, under-evaluated
starting models in seven papers, the modern equivalent of flaw F7. \textcite{shao2025spurious} show
that a randomly drawn reward improves Qwen models in mathematics, with no such gain for Llama or OLMo
models. \textcite{yue2025rlvr} show that the pass@$k$ measure reverses a ranking depending on $k$, an
instance of flaw F8.

\section{The vocabulary of language models}\label{sec:vocabulaire}

The vocabulary of RL has spread to the generative AI models that write text. Some of its algorithms
are borrowed, with only a loose resemblance to the original setting, to fine-tune the answers of a
pretrained language model, with no agent acting in a world. We fix our words as follows.

\begin{itemize}
  \item \textbf{Classical RL.} RL in the sense of Section~\ref{sec:cadre}: an agent, often trained from
  scratch, acts in an environment whose states evolve, in robotics, games or control. Offline RL
  belongs to it. Classical RL is the \file{classic} family of criterion X5.
  \item \textbf{RL on large models.} A label of convenience for the papers that call RL what they
  apply to a pretrained generative model conditioned on language (language model, vision-language
  model, text-to-image model). It is the \file{llm} family of criterion X5
  (Chapter~\ref{chap:03-sondage-methode}). The label takes up the papers' own word. It passes no
  judgement on the nature of what they do.
  \item \textbf{RLHF.} Learning from human feedback (\emph{reinforcement learning from human feedback})
  \parencite{christiano2017deep,ouyang2022instructgpt}. Annotators rank answers, a reward model learns
  their preferences, then the language model is fine-tuned to maximise this score. The acronym contains
  the letters RL. We nonetheless keep RLHF out of classical RL: the optimised policy produces a text in
  response to a prompt, then receives one score, with no environment whose state evolves under its
  actions.
  \item \textbf{RLVR and GRPO.} RLVR (\emph{reinforcement learning with verifiable rewards}) replaces
  the reward model with a program that checks the answer: the final answer to a mathematics problem
  compared with the correct one, code run against tests \parencite{deepseek2025r1}. GRPO (\emph{Group
  Relative Policy Optimization}) \parencite{shao2024deepseekmath} is the most widespread algorithm of
  this family. It derives from PPO: it draws a group of answers for each prompt and replaces the critic
  with the reward of each answer, standardised within the group. Chapter~\ref{chap:07-modeles-de-langue}
  describes our GRPO strand.
\end{itemize}

In this report, large language models are generative AI models. The readers of our survey are agents
built on such a model (Chapter~\ref{chap:03-sondage-methode}). There, the word \enquote{agent} denotes
a piece of software that reads a PDF and fills in a grid, unrelated to the agent of a Markov decision process.
